\documentclass{situationist-preprint}

\title{Free Will Under Indeterminism}
\msid{ACT-01-c}
\msfamily{ACT-01}
\msdiscipline{Action, Agency and Free Will}
\mstype{Alternative proof}
\mskeywords{free will; indeterminism; reasons-responsiveness; luck; rollback argument; deterministic chance; observational equivalence}

\begin{document}
\maketitle

\begin{abstract}
ACT-01-a derives the compatibility of free will with determinism and with indeterminism from a single premise, Independence, which asserts outright that worlds of both kinds contain agents with owned, reasons-responsive mechanisms. Critics reply that the premise restates the thesis. We reach the family's conclusion without it, with indeterminism in the foreground. Free action supervenes on the agent-level control profile of a mechanism together with its history; reasons-responsiveness is a threshold property of that profile, so a mechanism that meets it with a margin keeps it under small perturbations; and agent-level profiles can be realized, to any precision, by deterministic and by indeterministic fundamental dynamics, as results on the observational equivalence of deterministic and stochastic descriptions show. Given one actual free agent, both halves of the compatibility thesis follow, whichever kind our world is. The threat to free will is unreliability at the level of the agent, which both kinds of dynamics can produce or avoid. We apply the result to the rollback argument and to torn decisions.
\end{abstract}

\section{The premise avoided}\label{s:avoid}

ACT-01 holds that we have free will, and that free will is compatible both with determinism and with indeterminism; what threatens it is neither. Its principal manuscript defines free action as action from a mechanism of practical reasoning that is the agent's own and responsive to reasons (ACT-01-a, Definitions~2.2--2.4), and proves the compatibility thesis (ACT-01-a, Theorem~3.2) from one premise:

\begin{quote}
\emph{Independence} (ACT-01-a, Assumption~3.1). Whether an agent's mechanism is reasons-responsive and her own does not depend on whether her world is deterministic: there are deterministic worlds containing agents with such mechanisms, and indeterministic worlds containing agents with such mechanisms.
\end{quote}

The formalization of ACT-01-a shows that the theorem follows from Independence and from nothing else, and the principal manuscript draws the right lesson: the debate is located at Independence, whose first half the consequence argument attacks and whose second half the luck objection attacks. The proof of the theorem is a single step from Independence, and a critic may fairly say that Independence is the compatibility thesis restated as an existence claim. The principal manuscript gives a case for it, but the case is not a derivation: it argues that the counterfactuals that constitute reasons-responsiveness survive determinism, and that indeterminism need not disturb their reliability, without showing from more basic premises that worlds of each kind contain free agents.

The present manuscript derives Independence instead of assuming it. We use four premises. The first is that some actual agents act freely, which the family's affirmative answer asserts and the principal manuscript defends in its account of the reactive attitudes (ACT-01-a, Chapter~7). The second is that whether an agent acts freely supervenes on facts at the level of the agent: the profile of probabilities with which her mechanism issues in actions across situations that differ in their reasons, and the history by which the mechanism became hers. The third is that reasons-responsiveness, so understood, is a threshold property, which an agent who meets it with some margin continues to meet under small changes in her profile. The fourth is that every agent-level profile can be realized, to any precision, both by deterministic and by indeterministic fundamental dynamics. From these, both halves of Independence follow, and with them the compatibility thesis, whichever kind of world ours turns out to be.

The route differs from the principal one in two further respects. It does not rely on any particular semantics for counterfactuals under determinism, such as the local-miracle account of \citet{Lewis1981} that ACT-01-a uses in its case for the first half of Independence. And it treats indeterminism as the primary case. The principal manuscript addresses indeterminism briefly, by analogy with a coin-sorting machine whose parts are subject to quantum fluctuation, and defers the luck objection to Volume~III. Here the indeterministic half is derived first and in detail, and the analysis is then applied to the main arguments against it.

Section~\ref{s:history} reviews the debate about free will and indeterminism. Section~\ref{s:frame} states the framework and Section~\ref{s:arg} the argument. Section~\ref{s:loc} examines the three places at which indeterminism may enter deliberation, and Section~\ref{s:rollback} the rollback and luck arguments. Section~\ref{s:obj} considers objections; Section~\ref{s:compare} compares the two routes; Section~\ref{s:not} says what the argument does not establish.

\section{Indeterminism and free will: the debate}\label{s:history}

\subsection{The dilemma}

The thought that indeterminism is no friend of freedom is as old as the thought that determinism is its enemy. Hume held that liberty, opposed to necessity rather than to constraint, is the same thing as chance, and that chance is universally allowed to have no existence \citep[VIII]{Hume1748}. \citet{Hobart1934} argued that free will requires determination and is inconceivable without it: an action not determined by the agent's character and motives would not be hers. \citet{Ayer1954} gave the argument the form of a dilemma that has organized the debate since: either an action is caused, in which case it is necessitated, or it is not, in which case it is an accident; and in neither case is the agent responsible. ACT-01-d argues that the first of these positions, that free will requires determinism, is false (ACT-01-d, Proposition~1.2). The second horn of Ayer's dilemma is the subject of this manuscript.

\subsection{Libertarian responses}

Libertarians must answer the second horn, since they hold that free will requires indeterminism. Three answers have been developed. The first places indeterminism early in deliberation. \citet{Dennett1978} described a model, which he offered to libertarians, on which the considerations that come to mind are generated by an undetermined process and are then evaluated by a deterministic one; \citet{Mele1995} developed a \enquote{modest libertarianism} along these lines. Since the agent's evaluation of whatever considerations arise is reliable, the indeterminism does not make the outcome a matter of chance in the troubling sense. Critics have objected that such a view gives the agent no more control than a compatibilist account does, since the decisive step, the evaluation, is determined \citep{Clarke2003}.

The second answer places indeterminism in the decision itself, in the cases where it matters most. \citet{Kane1996} argued that in self-forming actions, which occur when an agent is torn between competing motivations, the agent makes an effort on behalf of each, and whichever effort succeeds, the outcome is something she was trying to bring about and for which she is responsible. \citet{Balaguer2010} defended a related view of \enquote{torn decisions}, arguing that whether they are undetermined is an open empirical question and that, if they are, they are free in the sense that matters. ACT-15 develops Kane's answer: indeterministic luck is compatible with responsibility when the agent tries to bring about each of the competing outcomes (ACT-15-a, Theorem~3.1), and the businesswoman of Kane's example satisfies its conditions (ACT-15-b, Proposition~2.2).

The third answer invokes agent causation: the agent, as a substance, causes her decision, and this is not reducible to causation by events in her. \citet{Clarke2003} examined the view in detail and concluded that it is coherent but that its advantages over event-causal libertarianism are less than its defenders claim. \citet{Franklin2018} has argued, in the other direction, that an event-causal libertarianism with minimal additions to compatibilist accounts of control is defensible against the luck objection.

\subsection{The luck objection and the rollback argument}

The most influential statement of the luck objection is the rollback argument of \citet{vanInwagen2000}. Suppose an agent faces a choice between telling the truth and lying, and that it is undetermined which she will do. Suppose God rolls the universe back to a moment before the choice a thousand times and lets it run forward. If the choice is undetermined, the outcomes will vary, and they will approach some stable ratio, say seventy truths to thirty lies. Watching the replays, we would be forced to conclude that what happened on each occasion was a matter of chance, and van Inwagen concluded that free will, which he held incompatible with determinism, is a mystery. \citet{Mele2006} developed the objection as a problem of \enquote{present luck}: nothing about the agent prior to the decision accounts for the difference between the world in which she decides one way and the world in which she decides the other, and so the difference is a matter of luck.

\citet{Pereboom2001,Pereboom2014} drew from these considerations the position of hard incompatibilism: free will is incompatible with determinism, and event-causal indeterminism does not help, since on such accounts the agent \enquote{disappears} at the moment of choice; nothing about her settles which option is taken. Only agent causation could supply the missing control, and we have no evidence that we are agent causes.

\subsection{Compatibilist treatments}

Compatibilists have mostly treated indeterminism as a side issue. \citet{FischerRavizza1998} note that their account of reasons-responsiveness is consistent with indeterminism, provided the indeterminism does not undermine the mechanism's reliability. The most developed recent treatment is that of \citet{List2014,List2019}, who argues that the possibilities relevant to free will are possibilities at the level of the agent, and that agent-level indeterminism, with genuine alternative possibilities, is compatible with physical determinism, because the agent-level state of a system does not fix the physical microstate and so does not fix what the system will do. The argument of the present manuscript is close to List's in its use of agent-level descriptions, but it runs in the opposite direction as well: it shows that agent-level profiles are equally compatible with physical indeterminism and physical determinism, and uses that fact to derive Independence.

\subsection{Determinism and its observational equivalents}

A final strand of the literature belongs to the philosophy of physics. Whether a theory is deterministic is a question about its laws \citep{Earman1986}. Whether the world is deterministic is not settled by the empirical success of a theory, since empirically equivalent theories may differ on it. The standard example is quantum mechanics: the collapse theory of \citet{GRW1986} is indeterministic, the pilot-wave theory of \citet{Bohm1952} is deterministic, and the two agree, at least to very high precision, on all observations performed so far. \citet{Suppes1993} argued that determinism is transcendental in the sense that no observations could decide between deterministic and indeterministic descriptions of certain systems, and \citet{Werndl2009} proved precise results of this kind: for a large class of deterministic systems there are stochastic processes that are observationally equivalent to them at every level of observational precision, and conversely. We use these results in Section~\ref{s:frame}.

\section{Framework}\label{s:frame}

\subsection{Agent-level profiles}

\begin{definition}[Reason-situation]\label{d:sit}
A \emph{reason-situation} for an agent is a specification of the considerations available to her that bear on what to do at a time. Two reason-situations differ if they differ in which reasons are present or in their strength.
\end{definition}

\begin{definition}[Control profile]\label{d:profile}
The \emph{control profile} of a mechanism of practical reasoning is the function $\pi$ that assigns to each reason-situation $s$ in a range $S$ and each action $a$ available in $s$ the probability $\pi(s,a)$ that the mechanism, operating as it actually operates, issues in $a$ when the agent is in $s$. The probability is fixed by the agent-level state of the agent in $s$: it is the measure, under the natural measure on the micro-realizations of that agent-level state, of the realizations that issue in $a$.
\end{definition}

Definition~\ref{d:profile} is the key move. The control profile is defined at the level of the agent: it concerns what her mechanism tends to do given how things are with her psychologically, not given the exact microstate of her brain. In an indeterministic world the probabilities are, or include, objective chances in the dynamics. In a deterministic world they are probabilities over micro-realizations, of the kind that statistical mechanics uses and that several authors have argued are objective chances in their own right \citep{Loewer2001,Glynn2010}. We do not need to settle whether such probabilities are chances. We need only that they are objective features of the agent-level state, fixed by the natural measure over its realizations, and the same in kind as the probabilities that cognitive science attributes to psychological processes.

\begin{definition}[Reasons-responsiveness with threshold]\label{d:rr}
Let $\tfrac12 < \theta \leq 1$. A mechanism with profile $\pi$ is \emph{$\theta$-reasons-responsive} if (a) in an understandable pattern of reason-situations in $S$ in which there is sufficient reason to do $a$, $\pi(s,a) \geq \theta$ (receptivity and reactivity to sufficient reasons); and (b) among these are situations in which there is sufficient reason to do otherwise than the agent actually does (reactivity to reasons to do otherwise). The mechanism meets the condition \emph{with margin} $\delta > 0$ if the inequalities in (a) hold with $\theta + \delta$ in place of $\theta$.
\end{definition}

Definition~\ref{d:rr} is the probabilistic form of Definition~2.2 of ACT-01-a, which follows \citet{FischerRavizza1998}: a mechanism is reasons-responsive if, operating as it actually operates, it would recognize sufficient reasons to act otherwise in a range of situations and would in some of them act on them. Where the profile takes only the values $0$ and $1$, Definition~\ref{d:rr} coincides with the principal definition for every $\theta$. Where it takes other values, \enquote{would act on them} is read as \enquote{would very probably act on them}, which is the reading on which an agent whose deliberation is subject to small agent-level noise still counts as responsive. We do not fix the value of $\theta$; the argument works for any value above one half. A value below one half would count as responsive a mechanism that more often than not ignored sufficient reasons, and no one has proposed that.

Two features of Definition~\ref{d:rr} matter for indeterminism. First, it concerns the \emph{pattern} of the profile across reason-situations, not its value in any one. A mechanism can assign probability one half to each of two options in a situation in which the reasons are balanced, and still be $\theta$-reasons-responsive, provided it acts on sufficient reasons where they are present. Second, condition (a) requires the pattern to be understandable: it must be such that a third party, knowing the agent's values, could see why she responds to these reasons and not to those \citep{FischerRavizza1998}. This is a constraint on the shape of the profile, not on its being deterministic.

\subsection{Principles and assumptions}

\begin{principle}[Agent-level supervenience]\label{p:sup}
Whether an agent acts freely on an occasion supervenes on (i) the control profile of the mechanism from which she acts, (ii) the history by which the mechanism came to be hers, described at the level of the agent, and (iii) the reason-situation she is in and the action she performs. Two agents alike in (i)--(iii) are alike in whether they act freely.
\end{principle}

\ref{p:sup} is implicit in the principal account. By Definition~2.4 of ACT-01-a, an agent acts freely if her action issues from a mechanism that is her own and reasons-responsive. Reasons-responsiveness is a property of the profile; ownership, by Definition~2.3, concerns whether the mechanism was installed or manipulated by another in a way that bypasses her own practical reasoning, and whether she takes responsibility for it, which are facts about her history at the level of the agent. \ref{p:sup} makes explicit that nothing else enters: in particular, not the deterministic or indeterministic character of the fundamental dynamics that realize the profile, except through the profile itself.

\begin{assumption}[Actuality with margin]\label{a:actual}
Some actual agent, on some occasion, acts from a mechanism that is her own and that is $\theta$-reasons-responsive with some margin $\delta > 0$.
\end{assumption}

The first part of Assumption~\ref{a:actual} is the family's affirmative answer, restricted to one case. Its support is the ordinary evidence on which we attribute agency to one another: people revise what they do when given reasons, in patterns that those who know them can understand, and much of their practical reasoning is not the product of manipulation. ACT-05 argues that the neuroscience of the readiness potential does not undermine this evidence. The second part, the margin, says only that the actual agent is not exactly at the threshold, so that the threshold is not met by luck at the last decimal place. Since $\theta$ is not fixed precisely, any agent who meets it comfortably meets it with a margin for some choice of $\theta$.

\begin{assumption}[Two-way realizability]\label{a:real}
For any world $w$ containing an agent $A$ in a reason-situation $s$ performing an action $a$, and any $\delta > 0$, there is a world $w'$ whose fundamental dynamics are indeterministic if those of $w$ are deterministic, and deterministic if those of $w$ are indeterministic, containing an agent $A'$ who has the same agent-level history as $A$, whose mechanism has a profile $\pi'$ with $|\pi'(s',b) - \pi(s',b)| < \delta$ for every reason-situation $s'$ and action $b$, and who is in $s$ and performs $a$.
\end{assumption}

Assumption~\ref{a:real} has two halves, and they rest on different considerations.

\emph{From determinism to indeterminism.} Take a deterministic world containing $A$. Add to its dynamics a stochastic perturbation whose magnitude is small enough that the probability, given any agent-level state, of each agent-level outcome changes by less than $\delta$. Agent-level processes are macroscopic and robust: they are realized by large numbers of neurons whose individual behavior is already noisy, and their agent-level outcomes are insensitive to sufficiently small microscopic perturbations. The result is an indeterministic world in which $A$'s counterpart has the same history and a profile within $\delta$ of hers. Since $\pi(s,a) > 0$ (she does perform $a$), and the perturbed profile is within $\delta$, we may choose $\delta$ small enough that $\pi'(s,a) > 0$, and among the indeterministic worlds so obtained there is one in which she performs $a$. The physical precedent is the collapse dynamics of \citet{GRW1986}, which adds to the deterministic Schr\"odinger evolution a stochastic process whose effects on macroscopic objects are, outside special circumstances, observationally negligible.

\emph{From indeterminism to determinism.} Take an indeterministic world containing $A$. The results of \citet{Werndl2009} show that for a large class of stochastic processes there are deterministic systems that are observationally equivalent to them at any given level of precision, and the agent-level state is a level of observation in the relevant sense: it is a coarse-graining of the microstate. The physical precedent is the pilot-wave theory of \citet{Bohm1952}, which reproduces the statistics of quantum mechanics from deterministic dynamics together with a probability distribution over initial configurations. In a deterministic world of this kind, the profile of $A$'s counterpart is given by the measure over micro-realizations of her agent-level state (Definition~\ref{d:profile}) and can match her original profile to any precision.

Assumption~\ref{a:real} is stronger than the physical results that support it, since those results concern particular classes of system and the assumption concerns all agent-level profiles. But the assumption concerns metaphysical possibility, not physical equivalence within a given theory, and the physical results show that nothing in the nature of agent-level probability ties it to one kind of fundamental dynamics. We return to the point in Section~\ref{s:obj}.

\section{The argument}\label{s:arg}

\begin{lemma}[Margin]\label{l:margin}
If a mechanism with profile $\pi$ is $\theta$-reasons-responsive with margin $\delta$, every mechanism with profile $\pi'$ such that $|\pi'(s,b) - \pi(s,b)| < \delta$ for every $s$ and $b$, and with the same pattern of sufficient reasons, is $\theta$-reasons-responsive.
\end{lemma}

\begin{proof}
For each reason-situation in the pattern of Definition~\ref{d:rr}(a), $\pi(s,a) \geq \theta + \delta$, so $\pi'(s,a) > \pi(s,a) - \delta \geq \theta$. The pattern of situations is unchanged, so it remains understandable, and the situations of Definition~\ref{d:rr}(b) remain among them.
\end{proof}

\begin{lemma}[Counterpart]\label{l:counter}
If $A$ acts freely in $w$ from a mechanism that is $\theta$-reasons-responsive with margin $\delta$, then for the world $w'$ and agent $A'$ given by Assumption~\ref{a:real} for $\delta$, $A'$ acts freely in $w'$.
\end{lemma}

\begin{proof}
By Assumption~\ref{a:real}, $A'$ has the same agent-level history as $A$, is in the same reason-situation and performs the same action, and her profile is within $\delta$ of $A$'s. By Lemma~\ref{l:margin}, her mechanism is $\theta$-reasons-responsive. Since her history is the same at the level of the agent, her mechanism is her own if $A$'s is. So $A'$ acts from an owned, reasons-responsive mechanism, and by \ref{p:sup} she is alike in freedom with $A$: she acts freely.
\end{proof}

\begin{theorem}[Indeterminism]\label{t:indet}
Given \ref{p:sup} and Assumptions~\ref{a:actual} and~\ref{a:real}, there is an indeterministic world containing an agent who acts freely. Free will is compatible with indeterminism.
\end{theorem}

\begin{proof}
By Assumption~\ref{a:actual}, some actual agent $A$ acts freely from a mechanism that is $\theta$-reasons-responsive with margin $\delta$. If the actual world is indeterministic, it is the required world. If it is deterministic, Assumption~\ref{a:real} gives an indeterministic world $w'$ containing a counterpart $A'$, who acts freely by Lemma~\ref{l:counter}.
\end{proof}

\begin{theorem}[Determinism]\label{t:det}
Given \ref{p:sup} and Assumptions~\ref{a:actual} and~\ref{a:real}, there is a deterministic world containing an agent who acts freely. Free will is compatible with determinism.
\end{theorem}

\begin{proof}
As for Theorem~\ref{t:indet}, with the roles of the two kinds of world exchanged.
\end{proof}

\begin{corollary}\label{c:indep}
Assumption~3.1 of ACT-01-a (Independence) holds, and with it Theorem~3.2 of ACT-01-a: free will is compatible both with determinism and with indeterminism.
\end{corollary}

\begin{proof}
Theorems~\ref{t:indet} and~\ref{t:det} give agents who act freely, and so have owned, reasons-responsive mechanisms, in worlds of each kind. That is the content of Independence, and the compatibility thesis follows from it by the proof of Theorem~3.2 of ACT-01-a.
\end{proof}

The proofs do not require us to know which kind of world is actual. This is the sense in which the result is robust to physics: whatever the fundamental dynamics turn out to be, the argument provides a world of the other kind in which a counterpart of an actual free agent is free. It also shows why Corollary~3.3 of ACT-01-a, that whether we have free will does not depend on whether determinism is true, is not merely a consequence of the compatibility thesis but an instance of a more general fact. The agent-level facts that settle freedom are invariant, up to arbitrarily small differences, under exchange of the fundamental dynamics.

\subsection{What threatens free will}

The family's conclusion includes a further claim: what threatens free will is neither determinism nor indeterminism. We now make this precise.

\begin{definition}[Threat]\label{d:threat}
A feature $F$ of worlds is a \emph{threat} to free will if, for some agent who acts freely in a world lacking $F$, the agent's counterpart, alike except for the presence of $F$ and what $F$ entails, does not act freely.
\end{definition}

\begin{proposition}\label{p:neither}
Neither determinism nor indeterminism of the fundamental dynamics is a threat to free will, except in so far as it entails a change in the agent's control profile or agent-level history. Every threat to free will operates through a change in the profile, the history, or the action.
\end{proposition}

\begin{proof}
By \ref{p:sup}, whether an agent acts freely is settled by (i)--(iii); so a feature that leaves these unchanged leaves freedom unchanged, and every threat must change one of them. By Assumption~\ref{a:real}, the presence of determinism or of indeterminism does not entail a change in (i) beyond any given $\delta$, nor any change in (ii) or (iii); and by Lemma~\ref{l:margin}, a change below the margin does not affect reasons-responsiveness.
\end{proof}

The threats that Proposition~\ref{p:neither} leaves are the familiar ones. Compulsion and addiction alter the profile, by making the mechanism insensitive to reasons that would otherwise move it; manipulation of the kind ACT-01-a discusses in Chapter~8 alters the history, by installing a mechanism in a way that bypasses the agent's own reasoning; and agent-level randomness, a profile that varies with something other than reasons, alters the profile by lowering its values below the threshold. The last is what the luck objection fears, and Proposition~\ref{p:neither} shows that it is not the same as indeterminism. A deterministic chaotic process in the brain, whose outcomes vary with microscopic differences unrelated to reasons, would threaten freedom in exactly the same way; and an indeterministic process confined to the generation of considerations, or to the resolution of genuinely balanced reasons, would not threaten it at all. What threatens free will is unreliability at the level of the agent, and unreliability is orthogonal to the distinction between determinism and indeterminism.

\section{Where indeterminism enters}\label{s:loc}

Indeterminism might enter the processes that realize deliberation in three places. We consider each in the light of Definition~\ref{d:rr}.

\subsection{Below the level of the agent}

If indeterminism is confined to microphysical processes whose agent-level effects are negligible, the agent's profile is the same, to within any margin, as it would be under deterministic dynamics. This is the case of the coin-sorting machine in ACT-01-a. Lemma~\ref{l:margin} shows that it does not affect freedom, and no libertarian or hard incompatibilist argument is directed at it, since it gives the agent no more alternatives at the agent level than determinism would.

\subsection{In the generation of considerations}

On the model of \citet{Dennett1978} and \citet{Mele1995}, which considerations come to the agent's mind is undetermined, while her evaluation of whatever considerations arise is reliable. On such a model, the probability that she acts on a sufficient reason, when one is present and comes to mind, is high; the probability that it comes to mind is less than one. Her profile, taken over reason-situations as Definition~\ref{d:sit} specifies them, assigns to each the probability that she acts on its reasons, and this is the product of the probability that they come to mind and the probability that, having come to mind, they move her.

\begin{proposition}\label{p:early}
On the early-indeterminism model, the agent's mechanism is $\theta$-reasons-responsive if and only if, in the relevant pattern of situations, the probability that a sufficient reason comes to mind, multiplied by the probability that it moves her once it has, is at least $\theta$.
\end{proposition}

\begin{proof}
Immediate from Definitions~\ref{d:profile} and~\ref{d:rr}.
\end{proof}

Proposition~\ref{p:early} shows that early indeterminism threatens freedom only if it makes the agent unreliable at noticing reasons, and that a deterministic agent who was equally unreliable at noticing them would be equally unfree. Attentive agents meet the condition. That reasons sometimes fail to come to mind is a fact about all human deliberation, under either kind of dynamics, and the threshold $\theta$ is a threshold for imperfect agents.

\subsection{In the decision}

On Kane's model, indeterminism enters at the point of decision, in cases where the agent is torn: she has sufficient reason for each of two options, and it is undetermined which she will choose. Her profile in such a situation $s$ assigns substantial probability to each option, say $\pi(s,a) = 0.6$ and $\pi(s,b) = 0.4$.

\begin{proposition}\label{p:torn}
An agent whose profile assigns intermediate probabilities to two options in torn situations, in which she has sufficient reason for each, can be $\theta$-reasons-responsive for every $\theta < 1$, provided that in the situations in which her reasons favor one option decisively she takes it with probability at least $\theta$.
\end{proposition}

\begin{proof}
Definition~\ref{d:rr} constrains the profile only in the pattern of situations in which there is sufficient reason for one action rather than another. In a torn situation there is sufficient reason for each option, and neither is the action that sufficient reasons require to the exclusion of the other. So the values of the profile in torn situations do not enter condition (a), and the condition is met if it is met in the situations where reasons are decisive.
\end{proof}

Proposition~\ref{p:torn} establishes that torn decisions are consistent with freedom in the sense of ACT-01: an agent who decides undeterminedly between two options, each of which she has sufficient reason to take, acts from a reasons-responsive mechanism, provided her mechanism is reliable where reasons are decisive. It does not by itself establish that she is responsible for the particular outcome. For that, ACT-15 adds the condition that she tries to bring about each of the competing outcomes, and shows that, when she does, she is responsible for whichever occurs (ACT-15-a, Theorem~3.1). The two results fit together. Freedom, in the sense of acting from an owned, reasons-responsive mechanism, is a property of her mechanism and is unaffected by the indeterminacy; responsibility for a particular outcome of a torn decision is secured by her effort towards each, which is the sense in which the outcome, whichever it is, was something she was trying to do. This division of labor is also the one that Volume~II of ACT-01 draws in its account of choice: on that account, having a choice about whether $p$ requires that $p$'s truth would vary with one's reasons through one's own mechanism (ACT-01, Volume~II, Chapter~14), and an outcome that would not vary with one's reasons, such as the landing of a coin that one can start but not influence, is not something about which one has a choice (ACT-01, Volume~II, Lemma~15.2, as corrected in ACT-01-e). The outcome of a torn decision is not of that kind: which way it goes varies with the agent's reasons across situations, even if not within the torn situation itself.

\section{The rollback and luck arguments}\label{s:rollback}

\subsection{Rollback}

Van Inwagen's rollback argument asks us to replay a single undetermined decision many times and observe the outcomes approach a stable ratio. In the terms of the present framework, the replays sample the profile at a single reason-situation $s$: the observed frequencies converge to $\pi(s,a)$ and $\pi(s,b)$.

\begin{proposition}\label{p:rollback}
The outcome of a rollback experiment on a single reason-situation is determined by the values of the agent's profile at that situation, and is consistent with any verdict about whether her mechanism is $\theta$-reasons-responsive.
\end{proposition}

\begin{proof}
The frequencies converge to $\pi(s,\cdot)$. By Definition~\ref{d:rr}, reasons-responsiveness depends on the values of $\pi$ across the pattern of situations in which reasons differ. For any values at $s$, there are profiles that agree with them at $s$ and are $\theta$-reasons-responsive (by Proposition~\ref{p:torn}, if $s$ is torn; by setting high values on the options sufficient reasons favor, if not), and profiles that agree with them at $s$ and are not (by setting low values elsewhere).
\end{proof}

Proposition~\ref{p:rollback} locates the error in the rollback argument. The argument invites us to conclude, from the variation of outcomes across replays of one situation, that the outcome is \enquote{just chance}. But the question whether an outcome is a matter of chance in the sense that excludes control is the question whether it varies with the agent's reasons, and this is a question about variation across situations with different reasons, which rollback holds fixed. A rollback of a deterministic agent would show the same outcome every time, and this would not show that she was free; a rollback of an indeterministic agent shows a ratio, and this does not show that she is unfree. The informative experiment is the one that varies the reasons: present the agent with stronger reasons for honesty, and see whether the ratio moves. If it moves in an understandable way, her mechanism is reasons-responsive, whatever happens in any single replay.

The point also clarifies what is right in van Inwagen's description. In a torn situation, a ratio of seventy to thirty is a fact about the agent: it records the relative weight that her reasons for each option have for her, and an agent for whom honesty weighed more would show a higher ratio. The stable ratio is therefore not evidence that her reasons play no role, but a measure of the role they play.

\subsection{Present luck}

\citet{Mele2006} asks us to compare two worlds with the same past up to the moment of decision, in one of which the agent decides to do $a$ and in the other $b$. Nothing about her before the decision explains why she decided $a$ in the first rather than $b$; the difference is a matter of luck.

The framework of this manuscript gives the reply that ACT-15-b gives for Kane's businesswoman: the demand for a contrastive explanation of the agent's deciding $a$ rather than $b$ in the same situation is not a condition of acting freely or of being responsible (ACT-15-b, abstract and Section~3). What explains her decision to do $a$ is her reasons for $a$, together with the reliability of her mechanism in acting on such reasons; what would have explained a decision to do $b$ is her reasons for $b$. There is no explanation of the contrast, because in a torn situation the contrast is not settled by anything prior; but each outcome is explained non-contrastively by reasons she had, and her profile across situations shows that her decisions vary with her reasons. By Proposition~\ref{p:neither}, the absence of contrastive explanation threatens freedom only if it lowers her profile below the threshold, which, in torn situations, it does not (Proposition~\ref{p:torn}).

\section{Objections and replies}\label{s:obj}

\begin{objection}[Begging the question]\label{o:beg}
\ref{p:sup} says that freedom supervenes on agent-level facts and is insensitive to the fundamental dynamics except through them. Incompatibilists deny exactly this: they hold that freedom depends on whether the past and the laws fix the action. The argument therefore assumes what Independence asserted, and the alternative proof has only moved the question.
\end{objection}

\begin{reply}
The objection is partly right, and it is worth being precise about how. Against the incompatibilist about determinism, \ref{p:sup} is contested, and the alternative proof does not dispense with the reply to the consequence argument. That reply is given in Volume~II of ACT-01, which shows that the argument's transfer rule, Rule Beta, is invalid (Theorem~16.3 there) and that repaired rules beg the question (ACT-01-b, Proposition~3.1). But the alternative proof changes what is in dispute. Independence is an existence claim about worlds of two kinds, which an incompatibilist can deny without denying anything else; \ref{p:sup} is a claim about what freedom consists in, which the incompatibilist must deny by identifying a fact, other than the agent's profile and history, on which freedom depends. For determinism, the candidate is the fixity of the action by the past and the laws, and the dispute is the one Volume~II addresses. For indeterminism, which is the topic of this manuscript, the candidate is less clear. The hard incompatibilist holds that an event-causal indeterministic agent lacks the power to settle which option occurs \citep{Pereboom2014}. Either settling is an agent-level fact, in which case it is a feature of the profile or the history and \ref{p:sup} accommodates it, or it is a fact beyond the agent level, such as agent causation by a substance, in which case the demand for it is a demand that the account of ACT-01 does not need to meet, since nothing in the definition of free action requires it. The proof therefore does more than move the question. It shows that the second half of Independence, which the luck objection attacks, follows from a supervenience principle that the luck objection does not itself deny.
\end{reply}

\begin{objection}[Observational equivalence is not agent-level equivalence]\label{o:obs}
The results of Werndl and Suppes concern observational equivalence: deterministic and stochastic descriptions agree on what can be observed at a given precision. The agent-level profile is not an observation. Two systems may agree on all observations while differing in their agent-level properties, so Assumption~\ref{a:real} does not follow from those results.
\end{objection}

\begin{reply}
The agent-level state is a coarse-graining of the microstate: it is fixed by which region of microstates the system is in, at a grain fixed by the psychological kinds. In the technical sense of \citet{Werndl2009}, an observation function is any function from states to a set of values that partitions the state space, and the map from microstates to agent-level states is such a function. The equivalence results therefore apply to it, for systems in the class they cover. The remaining gap, that the results cover particular classes of system and the assumption covers all profiles, is real; but it concerns how general the physical support is, not whether the support is of the right kind. The physical examples show that the agent-level probabilities of systems like ours can be realized by deterministic and by indeterministic dynamics, which is what the argument needs for agents like us. Assumption~\ref{a:real} generalizes this to every profile, and the generalization is required only for the agents to which Assumption~\ref{a:actual} applies, who are agents like us.
\end{reply}

\begin{objection}[Deterministic probabilities are not real]\label{o:detchance}
In a deterministic world, every action has probability one or zero given the complete microstate. The probabilities of Definition~\ref{d:profile} are therefore merely measures of our ignorance, and the profile of a deterministic agent is really a profile of zeros and ones. The two kinds of world do not share profiles.
\end{objection}

\begin{reply}
Definition~\ref{d:profile} defines the profile by the measure over micro-realizations of an agent-level state, and that measure is an objective feature of the state, not of anyone's ignorance: it is the same kind of quantity that statistical mechanics uses to explain why gases expand. Whether it is a chance is disputed \citep{Loewer2001,Glynn2010}, but freedom on the principal account is a matter of the agent's reasons-responsiveness, and the counterfactuals that constitute reasons-responsiveness are evaluated at the agent level. Even if one insists that a deterministic agent's \enquote{real} profile is a profile of zeros and ones, the zeros and ones are distributed over micro-realizations of her agent-level state in proportions given by the measure, and the agent-level profile summarizes them. A deterministic agent whose micro-realizations issue in the honest action in ninety-five per cent of the measure, whenever the reasons for honesty are decisive, is as reasons-responsive as an indeterministic agent with a chance of ninety-five per cent; and by \ref{p:sup} it is the agent-level fact that matters.
\end{reply}

\begin{objection}[Manipulated counterparts]\label{o:manip}
A counterpart with the same profile could have been produced by manipulation, and would not be free. So sameness of profile does not guarantee sameness of freedom, and Lemma~\ref{l:counter} fails.
\end{objection}

\begin{reply}
\ref{p:sup} includes the agent-level history among the supervenience base, and Assumption~\ref{a:real} requires the counterpart to share it. A manipulated agent differs in history: her mechanism was installed in a way that bypasses her own practical reasoning, which is the condition that Definition~2.3 of ACT-01-a excludes. Lemma~\ref{l:counter} applies only to counterparts with the same history, and these are not manipulated if the original was not. Manipulation is a threat to freedom in the sense of Definition~\ref{d:threat}, and Proposition~\ref{p:neither} classifies it correctly, as a threat that operates through the history.
\end{reply}

\begin{objection}[Ultimate responsibility]\label{o:ur}
Libertarians such as Kane hold that indeterminism is required not for reasons-responsiveness but for ultimate responsibility: for the agent to be the ultimate source of her character, some of her character-forming choices must be undetermined. The argument shows only that indeterminism does not undermine reasons-responsiveness, which libertarians never doubted.
\end{objection}

\begin{reply}
The objection concerns whether free will requires indeterminism, not whether it is compatible with it. ACT-08 argues that ultimate responsibility in the sense of the Basic Argument is impossible and is not required for responsibility, since self-formation can begin from an unchosen nature and proceed by degrees (ACT-08-a, Theorem~3.2 and Proposition~3.3). The present manuscript is consistent with the libertarian's view that indeterministic self-forming choices occur and contribute to self-formation; Propositions~\ref{p:torn} and~\ref{p:rollback} show that such choices can be free. It denies only that their indeterminism is what makes them free, which is the claim of Proposition~\ref{p:neither}.
\end{reply}

\section{The two routes compared}\label{s:compare}

The premises of the two routes can be set side by side.

\begin{center}\small
\begin{tabular}{@{}p{0.5\textwidth}cc@{}}
\toprule
Premise & ACT-01-a & ACT-01-c \\
\midrule
Definition of free action (ACT-01-a, Def.~2.4) & used & used \\
Independence (ACT-01-a, Assumption~3.1) & used & derived \\
Lewis's semantics for counterfactuals under determinism & used in support & not used \\
Agent-level supervenience (\ref{p:sup}) & implicit & used \\
Actuality with margin (Assumption~\ref{a:actual}) & not used & used \\
Two-way realizability (Assumption~\ref{a:real}) & not used & used \\
\bottomrule
\end{tabular}
\end{center}

The principal route proves the compatibility thesis without asserting that anyone is actually free; the actual freedom asserted by the family's answer is a further claim. The present route uses actual freedom as a premise, and in exchange derives Independence. This is not circular: Assumption~\ref{a:actual} asserts freedom of an agent in one world, whose kind we need not know, and the argument transfers it to worlds of both kinds. A critic who accepts that some of us are free, which every compatibilist and every libertarian does, and who accepts \ref{p:sup} and Assumption~\ref{a:real}, must accept the compatibility thesis whatever she believes about physics. In particular, a libertarian who believes that we are free, and that our world is indeterministic, must accept that free will is compatible with determinism, unless she denies \ref{p:sup} or Assumption~\ref{a:real}; and the denial of Assumption~\ref{a:real} would commit her to the view that agent-level profiles like ours cannot be realized by deterministic dynamics, which Bohmian mechanics, among other theories, makes implausible.

\section{What the argument does not establish}\label{s:not}

The argument does not establish that the actual world is deterministic or that it is indeterministic. It is designed to be neutral on the question, and Corollary~\ref{c:indep} holds either way.

It does not establish that indeterminism adds anything to freedom. A libertarian may hold that undetermined torn decisions are more fully the agent's own than determined ones \citep{Kane1996,Balaguer2010}. Proposition~\ref{p:neither} implies that, if this is so, the difference must show up in the agent's profile or history; it does not show that no such difference exists.

It does not establish responsibility for every outcome of an undetermined decision. Freedom, in the sense of acting from an owned, reasons-responsive mechanism, is compatible with indeterminism by Theorem~\ref{t:indet}; responsibility for a particular outcome of a torn decision requires in addition the dual effort of ACT-15 (ACT-15-a, Theorem~3.1). The full treatment of the luck objection is the subject of Volume~III of ACT-01.

Finally, it does not refute agent-causal libertarianism. If agents are substance causes, the argument is silent on whether that adds to their freedom. It shows only that the freedom ACT-01 attributes to us does not require it.

\section{Conclusion}\label{s:concl}

The compatibility thesis of ACT-01 can be derived without its principal premise. If freedom supervenes on the agent-level control profile and history of a mechanism, if reasons-responsiveness is a threshold property that a margin protects against small changes, if agent-level profiles can be realized both by deterministic and by indeterministic dynamics, and if some actual agent is free, then there are free agents in worlds of both kinds, whichever kind ours is. Indeterminism, located below the agent, in the generation of considerations or in torn decisions, is consistent with reasons-responsiveness; the rollback argument samples a profile at one situation and so cannot show its absence. What threatens free will is unreliability at the level of the agent, together with compulsion and manipulation, and none of these is determinism or indeterminism.

\begin{thebibliography}{99}\small
\bibitem[Ayer(1954)]{Ayer1954} Ayer, A.~J. (1954). Freedom and necessity. In \emph{Philosophical Essays} (pp.~271--284). Macmillan.
\bibitem[Balaguer(2010)]{Balaguer2010} Balaguer, M. (2010). \emph{Free Will as an Open Scientific Problem}. MIT Press.
\bibitem[Bohm(1952)]{Bohm1952} Bohm, D. (1952). A suggested interpretation of the quantum theory in terms of \enquote{hidden} variables, I. \emph{Physical Review}, 85(2), 166--179.
\bibitem[Clarke(2003)]{Clarke2003} Clarke, R. (2003). \emph{Libertarian Accounts of Free Will}. Oxford University Press.
\bibitem[Dennett(1978)]{Dennett1978} Dennett, D.~C. (1978). On giving libertarians what they say they want. In \emph{Brainstorms: Philosophical Essays on Mind and Psychology}. Bradford Books.
\bibitem[Earman(1986)]{Earman1986} Earman, J. (1986). \emph{A Primer on Determinism}. Reidel.
\bibitem[Fischer and Ravizza(1998)]{FischerRavizza1998} Fischer, J.~M. and Ravizza, M. (1998). \emph{Responsibility and Control: A Theory of Moral Responsibility}. Cambridge University Press.
\bibitem[Franklin(2018)]{Franklin2018} Franklin, C.~E. (2018). \emph{A Minimal Libertarianism: Free Will and the Promise of Reduction}. Oxford University Press.
\bibitem[Ghirardi et~al.(1986)]{GRW1986} Ghirardi, G.~C., Rimini, A. and Weber, T. (1986). Unified dynamics for microscopic and macroscopic systems. \emph{Physical Review D}, 34(2), 470--491.
\bibitem[Glynn(2010)]{Glynn2010} Glynn, L. (2010). Deterministic chance. \emph{The British Journal for the Philosophy of Science}, 61(1), 51--80.
\bibitem[Hobart(1934)]{Hobart1934} Hobart, R.~E. (1934). Free will as involving determination and inconceivable without it. \emph{Mind}, 43(169), 1--27.
\bibitem[Hume(1748)]{Hume1748} Hume, D. (1748). \emph{An Enquiry Concerning Human Understanding}, Section VIII.
\bibitem[Kane(1996)]{Kane1996} Kane, R. (1996). \emph{The Significance of Free Will}. Oxford University Press.
\bibitem[Lewis(1981)]{Lewis1981} Lewis, D. (1981). Are we free to break the laws? \emph{Theoria}, 47(3), 113--121.
\bibitem[List(2014)]{List2014} List, C. (2014). Free will, determinism, and the possibility of doing otherwise. \emph{No\^us}, 48(1), 156--178.
\bibitem[List(2019)]{List2019} List, C. (2019). \emph{Why Free Will Is Real}. Harvard University Press.
\bibitem[Loewer(2001)]{Loewer2001} Loewer, B. (2001). Determinism and chance. \emph{Studies in History and Philosophy of Modern Physics}, 32(4), 609--620.
\bibitem[Mele(1995)]{Mele1995} Mele, A.~R. (1995). \emph{Autonomous Agents: From Self-Control to Autonomy}. Oxford University Press.
\bibitem[Mele(2006)]{Mele2006} Mele, A.~R. (2006). \emph{Free Will and Luck}. Oxford University Press.
\bibitem[Pereboom(2001)]{Pereboom2001} Pereboom, D. (2001). \emph{Living Without Free Will}. Cambridge University Press.
\bibitem[Pereboom(2014)]{Pereboom2014} Pereboom, D. (2014). \emph{Free Will, Agency, and Meaning in Life}. Oxford University Press.
\bibitem[Suppes(1993)]{Suppes1993} Suppes, P. (1993). The transcendental character of determinism. \emph{Midwest Studies in Philosophy}, 18, 242--257.
\bibitem[van Inwagen(2000)]{vanInwagen2000} van Inwagen, P. (2000). Free will remains a mystery. \emph{Philosophical Perspectives}, 14, 1--19.
\bibitem[Werndl(2009)]{Werndl2009} Werndl, C. (2009). Are deterministic descriptions and indeterministic descriptions observationally equivalent? \emph{Studies in History and Philosophy of Modern Physics}, 40(3), 232--242.
\end{thebibliography}

\end{document}
